\cotefile{10}
\scanpage{185}

% slanted
\margin{$k$ \unsure{infini} ???}

Soit $k$ un corps. $k$ est un groupe algébrique sur $k$, dont nous cherchons les représentations linéaires, i.e les applications \uline{polynômes} $P(x)$ à valeurs dans un $\unsure{Gl|GL}(k^n)$, telles que $P(x+y)=P(x)P(y)$.

Les représentations \uline{irréductibles}. Alors le commutant \deleted{\unsure{des}} de l'ens. des $P(x)$ est un corps, \deleted{\unsure{soit donc|comme donc}} soit $D$, et les $P(x)$ prennent leurs valeurs dans le centre $\unsure{Z|\mathfrak z}$ de $D$. \inserted{(Comme $P$ est irréduc., on sait aussitôt que $D=Z$)}

$Z$ est un corps \inserted{\uline{commutatif}} fini sur $k$, donc on a une application $k$-polynômiale de $k$ dans une \deleted{\unsure{corps $Z$}}\inserted{\unsure{extension}} commutative $Z$ de $k$, telle que $P(x+y)=P(x)P(y)$. \deleted{\unsure{Si donc un|Soit donc un}} Je dis que $P(x)=1$ pour tt $x$ ;

\deleted{\unsure{étendant} les corps des scalaires à la clôture algébrique $\hat{k}$ de $k$, on obtient un $\hat{P}(x)$ hom. \inserted{\unsure{polynômial|polynôme}} de $\hat{k}$ dans le groupe multiplicatif des \unsure{élémts} inversibles de la $\hat{k}$-algèbre de dim. finie $A=\hat{k}\otimes_k Z$ \unsure{qui est comm. unitaire|On est ramené aussitôt} \unsure{finie}.}

\deleted{\unsure{On a}} \unsure{on met}

\[\begin{aligned} &\deleted{P(x)=a_n x^n+a_{n-1}x^{n-1}\unsure{\ldots+|+|+\ldots+}a_{\unsure{1|0}}x+1,\quad(\unsure{a_i|a}\in Z,\ x\in k)}\\ &\deleted{\text{avec }a_n\ne 0} \end{aligned}\]

\[\begin{aligned}&\deleted{P(x+y)=a_n x^n+a_n y^n\unsure{+\ldots}}\\ &\deleted{P(x)P(y)=a_n^2 x^n y^n+\unsure{P(x)|\ldots}}\end{aligned}\]

$P$ \unsure{serait} un polynôme de degré $n>0$, $P(x+y)$ \unsure{serait} de degré $n$ et $P(x)P(y)$ de degré $2n$, absurde.

% archive sheet number: 184
